%% ============================================================
%%  frontmatter/abstract.tex
%%  300 words or less - Single spacing - 11pt body text
%% ============================================================
\chapter*{Abstract}
\addcontentsline{toc}{chapter}{Abstract}
\thispagestyle{plain}

%% Abstract title: Title Case, 12pt Bold
{\fontsize{12}{15}\bfseries\selectfont
  Continuous Keystroke Authentication and Behaviour Analysis Using Keystroke Data for Programming Assessments\par}

\vspace{10pt}

%% Abstract body: 11pt, single spaced
\begin{singlespace}
\fontsize{11}{14}\selectfont

Online programming assessments need stronger identity assurance than password login while also preserving evidence of how a learner produced code. This dissertation presents BLAIM, a Behaviour Learning Analytics and Identity Monitoring module that performs continuous keystroke authentication and behaviour analysis using keystroke data captured during programming work. The system was implemented inside the GradeLoop Core microservice environment as a deployable FastAPI service with PostgreSQL persistence, Redis session buffering, TypeNet-style biometric embeddings, multi-phase enrollment, verification, identification, monitoring, archiving, and evaluator dashboards.

The module converts timestamped key events into temporal features including hold latency, inter-key latency, press latency, release latency, and normalized key code. A TypeNet-inspired LSTM model generates 128-dimensional embeddings from 70-event sequences. Enrollment stores a student template as the mean and standard deviation of multiple embeddings, while verification compares live sequence embeddings against the enrolled template using cosine similarity. The same event stream is also analysed for deletion rate, pause structure, paste/copy behaviour, burst typing, rhythm consistency, friction points, cognitive load indicators, authenticity, engagement, active problem solving, and learning depth.

The research follows a design science methodology: literature review, gap analysis, requirement derivation, architecture design, implementation, and evaluation using recorded real-user testing data. The evaluation dataset contains 32 enrollment records, 31 test records, 34 archived sessions, 12 authentication timeline events, and 16 distinct observed users. For labelled tests, the system produced 14 true accepts, 6 true rejects, 3 false rejects, and 1 false accept at a 0.70 threshold, giving an 82.35\% genuine accept rate and a 14.29\% false positive rate on the impostor tests.

The main contribution is an integrated, interpretable, and LMS-ready approach that combines biometric identity evidence with learning-process evidence from one keystroke stream. Deployment considerations include device sensitivity, threshold calibration, and formal fairness governance before institution-wide use.

\bigskip

\noindent\textbf{Keywords:} keystroke dynamics, continuous authentication, behaviour analysis, behavioural biometrics, programming education

\end{singlespace}
\clearpage
